\chapter{Example Circuit and Testing}

\section{Overview}

Two example circuits were used to validate the complete OpenROAD integration workflow: \textbf{FullAdder} and \textbf{HalfAdder}. Both circuits were successfully processed through the complete RTL-to-GDSII backend flow and the generated outputs were verified at each implementation stage.

\section{FullAdder Testing}

\subsection{Input: FullAdder.cir.out}

The FullAdder project uses a \texttt{.cir.out} netlist generated from the KiCad schematic inside eSim. The circuit implements a 1-bit FullAdder with three input ports (\texttt{in1}, \texttt{in2}, and \texttt{cin}) and two output ports (\texttt{sum} and \texttt{cout}).

\subsection{Generated Verilog}

The \texttt{netlist2rtl.py} converter processed the generated \texttt{FullAdder.cir.out} file and automatically generated the following Verilog module:

\begin{lstlisting}[language=verilog,
caption={Generated FullAdder.v module.}]
module FullAdder (
    cin,
    cout,
    in1,
    in2,
    sum
);

input cin;
output cout;
input in1;
input in2;
output sum;

wire net1;
wire net2;
wire net3;
wire net4;
wire net5;

sky130_fd_sc_hd__fa_1 _0_ (
    .A(net2),
    .B(net3),
    .CIN(net1),
    .COUT(net4),
    .SUM(net5)
);

sky130_fd_sc_hd__clkdlybuf4s50_1 input1 (
    .A(cin),
    .X(net1)
);

sky130_fd_sc_hd__clkdlybuf4s50_1 input2 (
    .A(in1),
    .X(net2)
);

sky130_fd_sc_hd__clkdlybuf4s50_1 input3 (
    .A(in2),
    .X(net3)
);

sky130_fd_sc_hd__clkdlybuf4s50_1 output4 (
    .A(net4),
    .X(cout)
);

sky130_fd_sc_hd__clkdlybuf4s50_1 output5 (
    .A(net5),
    .X(sum)
);

endmodule
\end{lstlisting}

\subsection{Generated SDC File}

The timing constraints file was automatically generated during ORFS setup.

\begin{lstlisting}[language=tcl,
caption={Auto-generated FullAdder constraint.sdc file.}]
current_design FullAdder

create_clock -name clk -period 10.0000

set_input_delay 0.1000 -add_delay \
[get_ports {cin}]

set_input_delay 0.1000 -add_delay \
[get_ports {in1}]

set_input_delay 0.1000 -add_delay \
[get_ports {in2}]

set_output_delay 0.1000 -add_delay \
[get_ports {cout}]

set_output_delay 0.1000 -add_delay \
[get_ports {sum}]
\end{lstlisting}

\subsection{Flow Completion}

The OpenROAD Flow Scripts execution for the FullAdder circuit completed successfully through all backend stages including synthesis, floorplanning, placement, CTS, routing, and GDSII generation.

Figure~\ref{fig:fa_flow_done} shows the terminal output after successful completion of the RTL-to-GDSII flow.

\begin{figure}[H]
\centering
\includegraphics[width=0.92\textwidth]{FA-Flow-completed.png}
\caption{Successful RTL-to-GDSII flow completion for the FullAdder circuit.}
\label{fig:fa_flow_done}
\end{figure}

\subsection{Generated Output Files}

After successful ORFS execution, multiple backend design files were generated automatically.

\begin{table}[H]
\centering
\caption{Generated output files after ORFS execution.}

\begin{tabular}{lll}
\toprule

\textbf{File} &
\textbf{Format} &
\textbf{Description} \\

\midrule

\texttt{FullAdder.v}
& Verilog
& Generated RTL module \\

\texttt{FullAdder.def}
& DEF
& Placement and routing data \\

\texttt{FullAdder.gds}
& GDSII
& Final layout file \\

\texttt{FullAdder.sdc}
& SDC
& Timing constraints \\

\texttt{FullAdder.spef}
& SPEF
& Parasitic extraction output \\

\texttt{logs/}
& Directory
& Backend execution logs \\

\bottomrule
\end{tabular}

\label{tab:output_files}
\end{table}

\section{OpenROAD GUI Visualization}

After completion of the backend flow, the generated layout was opened automatically in the OpenROAD GUI using the generated \texttt{.odb} database file.

Figure~\ref{fig:fa_openroad} shows the FullAdder physical layout displayed in the OpenROAD GUI.

\begin{figure}[H]
\centering
\includegraphics[width=0.92\textwidth]{FA-OpenROAD.png}
\caption{FullAdder layout visualized in the OpenROAD GUI.}
\label{fig:fa_openroad}
\end{figure}

\section{KLayout Visualization}

The generated GDSII layout file was also verified using KLayout.

Figure~\ref{fig:fa_klayout} shows the FullAdder layout loaded in KLayout for physical inspection.

\begin{figure}[H]
\centering
\includegraphics[width=0.92\textwidth]{FA-Klayout.png}
\caption{FullAdder GDS layout visualized in KLayout.}
\label{fig:fa_klayout}
\end{figure}

\section{HalfAdder Testing}

The HalfAdder circuit was also validated successfully using the same OpenROAD integration workflow. The generated \texttt{HalfAdder.cir.out} file was converted into Verilog using \texttt{netlist2rtl.py}, and the OpenROAD Flow Scripts execution completed successfully.

The generated backend outputs included:

\begin{itemize}
    \item synthesized Verilog,
    \item DEF layout files,
    \item GDSII layout files,
    \item timing constraint files,
    \item ORFS execution logs.
\end{itemize}

Figure~\ref{fig:fa_ha} shows both FullAdder and HalfAdder projects after successful flow completion.

\begin{figure}[H]
\centering
\includegraphics[width=0.92\textwidth]{FA-HA.png}
\caption{FullAdder and HalfAdder after successful RTL-to-GDSII flow execution.}
\label{fig:fa_ha}
\end{figure}

\section{Workflow Summary}

The complete validated workflow implemented during the internship is summarized below:

\begin{enumerate}
    \item Create and simulate the circuit in eSim using KiCad and Ngspice.
    
    \item Run KiCad-to-Ngspice conversion to generate the \texttt{.cir.out} file.
    
    \item Click the OpenROAD button integrated into the eSim toolbar.
    
    \item Automatically generate Verilog RTL using \texttt{netlist2rtl.py}.
    
    \item Automatically generate \texttt{config.mk} and \texttt{constraint.sdc}.
    
    \item Execute the ORFS backend implementation flow.
    
    \item Generate backend outputs including DEF and GDSII layouts.
    
    \item Visualize the generated layouts using OpenROAD GUI and KLayout.
\end{enumerate}